\documentclass[12pt]{article}
\usepackage[utf8]{inputenc}
\usepackage[T1]{fontenc}
\usepackage{hyperref}
\title{Kamas Language Resources for K-12 Curriculum Development}
\author{Compiled by Claude Sonnet 5.5 (claude-sonnet-5-5) for Andr\'as Kornai}
\date{}
\begin{document}
\maketitle

\section{Introduction}
Kamas (Kamassian; ISO 639-3 xas; Glottocode kama1351) is an extinct Southern Samoyedic language of the Uralic family, formerly spoken in Russia (Siberia). Its last known speaker, Klavdiya Plotnikova, died in 1989. It is better documented than most extinct languages, through Kai Donner's fieldwork (1912--1914) and recordings of the last speaker (1964--1970). No school programme or revitalization curriculum was found.

\section{Primary Educational Resources}
No primers, textbooks or classroom materials for Kamas were found. A teaching-oriented PDF from LMU Munich's e-learning site (``o1 kamas'') appeared in the results; I did not verify its contents. \url{https://www.infuse.finnougristik.uni-muenchen.de/e-learning/kamas/o1_kamas.pdf}

\section{Linguistic and Reference Materials}
\begin{itemize}
\item Wikipedia, ``Kamas language''. \url{https://en.wikipedia.org/wiki/Kamas_language}
\item Kai Donner's \emph{Kamassisches W\"orterbuch} (ed.\ Joki 1944, Lexica Societatis Fenno-Ugricae 8), known to me from a bibliographic listing at ASJP, not as an online text. \url{https://asjp.clld.org/languages/KAMAS}
\item Documentation of the INEL Kamas corpus (version 2.0). \url{https://s3-uhh.lzs.uni-hamburg.de/gwiss-inel-corpora/remote/kamas-2.0/documentation/kamas-2.0-documentation.pdf}
\item Article on Kamas from journal.fi (Finnish journal; exact title not confirmed). \url{https://journal.fi/fuf/article/download/113108/67507/222194}
\item Endangered Languages Project entry for Kamas. \url{https://www.endangeredlanguages.com/node/109155}
\end{itemize}

\section{Religious and Community Resources}
No scripture or community language organization was found.

\section{Digital and Audio Resources}
The INEL Kamas corpus has two parts, Donner's folklore texts and transcribed audio recordings of the last speaker from 1964--1970. It is available under a CC-BY-NC-SA licence. \url{https://inel.corpora.uni-hamburg.de/KamasCorpus}

\section{Curriculum Development Context}
The INEL corpus gives Kamas a documentary base that could support a heritage unit with texts and recordings. No native speakers remain, so instruction in the language is not possible beyond what the corpus and dictionary allow. Descendant communities (Kamassians) are the natural partners, but none were identified in these searches.
\end{document}